% Part: first-order-logic
% Chapter: axiomatic-deduction
% Section: axioms-rules-propositional

\documentclass[../../../include/open-logic-section]{subfiles}

\begin{document}

\iftag{FOL}
      {\olfileid{fol}{axd}{prp}}
      {\olfileid{pl}{axd}{prp}}

\olsection{ప్రతిజ్ఞాత్మక సంయోజకాల స్వీకృతాలు మరియు నియమాలు}

\begin{defn}[స్వీకృతాలు]
ప్రతిజ్ఞాత్మక సంయోజకాల \emph{స్వీకృతాల} సమితి~$\PAx$ కింది
రూపాలన్నింటిలోని !!{formula}sను కలిగి ఉంటుంది:
\begin{align}
  & (!A \land !B) \lif !A \ollabel{ax:land1}\\
  & (!A \land !B) \lif !B \ollabel{ax:land2}\\
  & !A \lif (!B \lif (!A \land !B)) \ollabel{ax:land3}\\
  & !A \lif (!A \lor !B) \ollabel{ax:lor1}\\
  & !A \lif (!B \lor !A) \ollabel{ax:lor2}\\
  & (!A \lif !C) \lif ((!B \lif !C) \lif ((!A \lor !B) \lif !C)) \ollabel{ax:lor3}\\
  & !A \lif (!B \lif !A) \ollabel{ax:lif1}\\
  & (!A \lif (!B \lif !C)) \lif ((!A \lif !B) \lif (!A \lif !C)) \ollabel{ax:lif2}\\
  & (!A \lif !B) \lif ((!A \lif \lnot !B) \lif \lnot !A) \ollabel{ax:lnot1}\\
  & \lnot !A \lif (!A \lif !B) \ollabel{ax:lnot2}\\
  & \ltrue \ollabel{ax:ltrue}\\
  & \lfalse \lif !A \ollabel{ax:lfalse1}\\
  & (!A \lif \lfalse) \lif \lnot !A \ollabel{ax:lfalse2}\\
  & \lnot\lnot !A \lif !A \ollabel{ax:dne}
\end{align}
\end{defn}

\begin{defn}[మోడస్ పోనెన్స్]
  ఒక !!a{derivation}లో $!B$, $!B \lif !A$ ఇప్పటికే కనిపిస్తే,
  $!A$ ఒక సరైన నిగమన దశ.
\end{defn}

మోడస్ పోనెన్స్ నియమాన్ని ``\MP'' అని సంక్షిప్తంగా రాస్తాం.

\end{document}
